% ===== PRÉAMBULE CANONIQUE — à coller VERBATIM en tête de CHAQUE .tex =====
\documentclass[11pt,a4paper]{article}
\usepackage[utf8]{inputenc}\usepackage[T1]{fontenc}\usepackage{lmodern}
\usepackage[french]{babel}
\usepackage{amsmath,amssymb,mathtools}\usepackage{icomma}
\usepackage{siunitx}\sisetup{locale=FR}  % \qty{}{}, \si{}, \ang{}
\usepackage{xcolor}\usepackage{colortbl}
\usepackage{tikz}\usepackage{pgfplots}\pgfplotsset{compat=1.18}
\usepackage[siunitx,europeanresistors]{circuitikz} % circuits (élec.)
\usepackage[most]{tcolorbox}\usepackage{enumitem}\usepackage{array,booktabs}
\usepackage[a4paper,margin=2.2cm]{geometry}
\usetikzlibrary{arrows.meta,calc,angles,quotes,positioning,patterns,%
  backgrounds,shapes.geometric,decorations.pathmorphing,decorations.markings,intersections}
\definecolor{primary}{RGB}{222,49,99}\definecolor{primarydark}{RGB}{150,22,55}
\definecolor{defcol}{RGB}{0,114,178}\definecolor{methcol}{RGB}{52,92,148}
\definecolor{excol}{RGB}{0,135,90}\definecolor{attncol}{RGB}{200,40,55}
\tcbset{boxbase/.style={enhanced,breakable,boxrule=0.9pt,arc=2.5pt,
  left=3mm,right=3mm,top=2.3mm,bottom=2.3mm,fonttitle=\bfseries,coltitle=white}}
% --- boîtes TUTORIEL (noms EXACTS, ne pas renommer) ---
\newtcolorbox{cle}[1][]{boxbase,colback=defcol!6,colframe=defcol,
  title={\textsc{À retenir}\ifx\relax#1\relax\relax\else\textmd{ : \itshape #1}\fi}}
\newtcolorbox{methode}[1][]{boxbase,colback=methcol!7,colframe=methcol,sharp corners,
  title={\textsc{Méthode}\ifx\relax#1\relax\relax\else\textmd{ --- \itshape #1}\fi}}
\newtcolorbox{exemplebox}[1][]{boxbase,colback=excol!5,colframe=excol,
  title={\textsc{Exemple}\ifx\relax#1\relax\relax\else\textmd{ : \itshape #1}\fi}}
\newtcolorbox{piege}[1][]{boxbase,colback=attncol!6,colframe=attncol,
  title={\textsc{Piège}\ifx\relax#1\relax\relax\else\textmd{ : \itshape #1}\fi}}
\newtcolorbox{remarque}[1][]{boxbase,colback=black!4,colframe=black!45,
  title={\textsc{Remarque}\ifx\relax#1\relax\relax\else\textmd{ : \itshape #1}\fi}}
% --- boîtes EXERCICE / CORRIGÉ (noms EXACTS, auto-comptées) ---
\newtcolorbox[auto counter]{exo}[1][]{boxbase,colback=primary!4,colframe=primary,
  title={\textsc{Exercice}~\thetcbcounter\ifx\relax#1\relax\relax\else\textmd{ --- \itshape #1}\fi}}
\newtcolorbox[auto counter]{corrige}[1][]{boxbase,colback=excol!4,colframe=excol,
  title={\textsc{Corrigé}~\thetcbcounter\ifx\relax#1\relax\relax\else\textmd{ --- \itshape #1}\fi}}
\newcommand{\vv}[1]{\vec{#1}}\newcommand{\ds}{\displaystyle}
\newcommand{\R}{\mathbb{R}}\newcommand{\N}{\mathbb{N}}\newcommand{\Z}{\mathbb{Z}}
% =========================================================================
\begin{document}
\usetikzlibrary{babel}
\begin{exo}[intensité à partir de la charge]
Une charge $Q=\qty{12}{\coulomb}$ traverse la section d'un fil conducteur en une durée
$t=\qty{4}{\second}$. Calcule l'intensité $I$ du courant, en précisant l'unité.
\end{exo}

\begin{exo}[de la charge au nombre d'électrons]
Un conducteur est parcouru par un courant continu d'intensité $I=\qty{2}{\ampere}$ pendant
$t=\qty{5}{\second}$. On donne la charge élémentaire $e=\qty{1.6e-19}{\coulomb}$.
\begin{enumerate}[label=\alph*)]
  \item Quelle charge $Q$ traverse une section du conducteur ?
  \item À combien d'électrons cette charge correspond-elle ?
\end{enumerate}
\end{exo}

\begin{exo}[loi des nœuds : une intensité manquante]
Au nœud $N$ ci-dessous arrive un courant d'intensité $I=\qty{5}{\ampere}$. Deux branches en partent :
l'une porte $I_1=\qty{2}{\ampere}$. Détermine $I_2$.
\begin{center}
\begin{circuitikz}[scale=0.9,transform shape]
  \coordinate (N) at (0,0);
  \draw[-{Stealth[length=2mm]},primary,line width=1pt] (-1.8,0) -- (-0.1,0)
     node[midway,above,font=\footnotesize,primary]{$I$};
  \draw[thick] (N) -- (1.4,0.9); \draw[-{Stealth[length=1.8mm]},primary] (0.5,0.32) -- (1.0,0.64)
     node[right,font=\footnotesize,primary]{$I_1$};
  \draw[thick] (N) -- (1.4,-0.9); \draw[-{Stealth[length=1.8mm]},primary] (0.5,-0.32) -- (1.0,-0.64)
     node[right,font=\footnotesize,primary]{$I_2$};
  \fill (N) circle(2pt); \node[above left,font=\footnotesize] at (N){$N$};
\end{circuitikz}
\end{center}
\end{exo}

\begin{exo}[en série, l'intensité est partout la même]
Dans le circuit série ci-dessous, un ampèremètre placé avant $R_1$ indique $\qty{0.5}{\ampere}$.
Que lira un second ampèremètre placé entre $R_1$ et $R_2$ ? Justifie en une phrase.
\begin{center}
\begin{circuitikz}[scale=0.9,transform shape]
  \draw (0,0) to[battery1,l_=$G$] (0,2)
        to[ammeter](1.4,2)
        to[R,l=$R_1$](3.2,2)
        to[R,l=$R_2$](5,2) -- (5,0) -- (0,0);
\end{circuitikz}
\end{center}
\end{exo}

\begin{exo}[isoler la durée]
Un courant continu d'intensité $I=\qty{0.5}{\ampere}$ circule dans un fil. Quelle durée $t$ faut-il
pour qu'une charge $Q=\qty{30}{\coulomb}$ traverse une section de ce fil ?
\end{exo}

\end{document}
